%% ---------------------------------------------------------------------------
\section{The Left End}\label{sec:left}

Near the left border, the reflected triangle is produced by the anti-diagonals
$t+i=2n-1$, $t+i=2n$, $\dots$, which follow the echo of the right end; we call
them \emph{chains}. The first few chains form a finite automaton that reads the
two input anti-diagonals $t+i=2n-3$ and $t+i=2n-2$ and must fire the cells near
the left border exactly when the echo reaches them. If the input anti-diagonals were periodic over a long
stretch, this automaton would repeat a state and fire too early
(\cref{lem:left}). Applied to all lengths at once, this excludes periodic
wedges at the left border of the half-line (\cref{thm:left}).

\subsection{Chains and Slices}

Fix $n\ge2$. For $c\ge0$ and $\max(1,c)\le i\le n$, and for $c=-1$ and $1\le
i\le n-1$, put
\[
v_c(i):=C_n(2n-2+c-i,\,i),
\]
the cell at position $i$ of the \emph{chain} $t+i=2n-2+c$. Chains $-1$ and $0$
are the input anti-diagonals, and chain $1$ is the return chain
(\cref{fig:chains}). The predecessors of the cell $(2n-2+c-i,i)$ lie on the
chains $c-2$, $c-1$ and $c$, so
\begin{equation}\label{eq:chains}
v_c(i)=\delta\bigl(v_{c-2}(i-1),\,v_{c-1}(i),\,v_c(i+1)\bigr)\qquad(c\ge1),
\end{equation}
where $v_{c-2}(0)$ and $v_c(n+1)$ stand for the borders. If the line of length
$n$ is synchronized at time $2n-2$, then the cell $i$ fires on chain $i$:
$v_i(i)=\Fq$, and $v_c(i)\ne\Fq$ for $c<i$. In the skewed coordinates
$u_c(y):=v_c(y+\lfloor c/2\rfloor)$ the recursion \eqref{eq:chains} becomes
\begin{equation}\label{eq:slices}
u_c(y)=\delta\bigl(u_{c-2}(y),\,u_{c-1}(y+\varepsilon_c),\,u_c(y+1)\bigr),
\qquad \varepsilon_c=1\text{ for even }c,\ \ \varepsilon_c=0\text{ for odd }c,
\end{equation}
with $u_0(y)=C_\infty(2n-2-y,y)$ and $u_{-1}(y)=C_\infty(2n-2-y,y-1)$. Here
$u_c(y)$ is the cell $y+\lfloor c/2\rfloor$ at time $2n-2+\lceil c/2\rceil-y$;
it is $\Fq$ if $\lceil c/2\rceil=y$ and not $\Fq$ if $\lceil c/2\rceil<y$. The
\emph{slice} $y$ is the sequence $(u_1(y),u_2(y),\dots)$. Its cells form a
staircase that starts just below the input cells $u_{-1}(y)$ and $u_0(y)$, at
the cell $y$ of chain $1$, and descends to the right: from $u_c(y)$ to
$u_{c+1}(y)$ the position increases by one if $c$ is odd, and the time
increases by one if $c$ is even. By \eqref{eq:slices}, each slice is computed
from the slice $y+1$ and the two input cells $u_{-1}(y)$ and $u_0(y)$: reading
the input anti-diagonals from right to left, the slices form a finite
automaton. A word $w$ on an
integer interval $I$ is \emph{$\ell$-periodic} if $w(y+\ell)=w(y)$ whenever
$y,y+\ell\in I$.

\begin{figure}[t]
\centering
%% Figure: the lower left corner of the line of length n = 12, coloured by chain, with the slices of
%% the proof of the left-border lemma for s = 2.
\begin{tikzpicture}
  \def\u{0.44}\def\v{0.33}
  \def\tz{13}
  % cell (t,i) occupies [(i-1)u, iu] x [-(t-tz+1)v, -(t-tz)v]
  \foreach \t in {13,...,22} {%
    \foreach \i in {1,...,10} {%
      \pgfmathtruncatemacro{\c}{\t+\i-22}
      \pgfmathtruncatemacro{\odd}{mod(\c+10,2)}
      \def\cc{rgAgree}%
      \ifnum\c>-2 \def\cc{rgInput}\fi
      \ifnum\c=1 \def\cc{rgReturn}\fi
      \ifnum\c>1 \ifnum\odd=1 \def\cc{rgRefl2}\else\def\cc{rgRefl}\fi\fi
      \ifnum\t=22 \def\cc{stF}\fi
      \filldraw[fill=\cc,draw=grid!45,line width=0.15pt]
        (\i*\u-\u,-\t*\v+\tz*\v) rectangle (\i*\u,-\t*\v+\tz*\v-\v);
    }}
  \draw[ink,line width=1.3pt] (0,0.06) -- (0,-10*\v-0.06);
  \def\cellbox#1#2#3{\draw[#3] (#2*\u-\u,-#1*\v+\tz*\v) rectangle (#2*\u,-#1*\v+\tz*\v-\v);}
  % slices y = 4 (dashed), y = 3 (thin) and y = s = 2 (thick), each with its two input cells
  \foreach \t/\i in {18/3,18/4,19/4,19/5,20/5,20/6} {\cellbox{\t}{\i}{ink,line width=0.7pt,densely dashed}}
  \foreach \t/\i in {19/2,19/3,20/3,20/4,21/4,21/5} {\cellbox{\t}{\i}{ink,line width=0.8pt}}
  \foreach \t/\i in {20/1,20/2,21/2,21/3,22/3,22/4} {\cellbox{\t}{\i}{ink,line width=1.8pt}}
  % on the firing row, cell i lies on chain i
  \foreach \i in {1,...,10} {\node[font=\scriptsize,text=white] at (\i*\u-0.5*\u,-9.5*\v) {\i};}
  % chains -1, 0, 1 cross the first row shown at the cells 8, 9, 10
  \node[font=\scriptsize,anchor=south] at (7.5*\u,0.03) {$-1$};
  \node[font=\scriptsize,anchor=south] at (8.5*\u,0.03) {$0$};
  \node[font=\scriptsize,anchor=south] at (9.5*\u,0.03) {$1$};
  \node[font=\scriptsize,anchor=south east] at (6.95*\u,0.03) {chain};
  \node[font=\scriptsize,anchor=east] at (-0.08,-9.5*\v) {$t=2n-2$};
  \node[font=\scriptsize,anchor=east] at (-0.08,-0.5*\v) {$t=13$};
  % legend
  \begin{scope}[font=\scriptsize]
    \node[anchor=west,inner sep=0pt] at (4.8,-0.15) {chain $c$: the anti-diagonal $t+i=2n-2+c$};
    \foreach \c/\l [count=\k from 1] in {
        rgInput/{chains $-1$ and $0$: the input anti-diagonals},
        rgReturn/{chain $1$: the return chain},
        stF/{firing row: cell $i$ fires on chain $i$}}
      \node[anchor=west,inner sep=0pt] at (4.8,-0.15-0.45*\k) {\swatch{\c}\ \l};
    \node[anchor=west,inner sep=0pt] at (4.8,-1.95) {\swatch{rgRefl}\,\swatch{rgRefl2}\ chains $c\ge2$, alternately shaded};
    \node[anchor=west,inner sep=0pt] at (4.8,-2.5) {outlined: slices $y=2$ (thick), $3$ and $4$ (dashed)};
  \end{scope}
\end{tikzpicture}

\caption{The lower left corner of the line of length $n=12$, coloured by
chain. The cell $i$ of the firing row lies on chain $i$. Each outlined group of
six cells is a slice $(u_1(y),\dots,u_4(y))$ of the band of chains
$1,\dots,4$ used in the proof of \cref{lem:left} with $s=2$, together with its
two input cells $u_{-1}(y)$ and $u_0(y)$: the slice is computed from the slice
to its right and from these two input cells. The group drawn with a thick line
is the one at $y=s$; its last cell $u_{2s}(s)$ is in the firing row.}
\label{fig:chains}
\end{figure}

\subsection{The Left-Border Lemma}

\begin{lemma}[left border]\label{lem:left}
Let $\delta$ have $k$ states and synchronize the line of length $n$ at time
$2n-2$. Let $s,\ell\ge1$ and $s\le Y\le n-s$, and suppose that the words
$y\mapsto C_\infty(2n-2-y,y)$ for $s\le y\le Y$ and $y\mapsto
C_\infty(2n-3-y,y)$ for $s\le y\le Y-1$ are $\ell$-periodic. Then
$Y-s\le(k-1)^{2s}\ell$.
\end{lemma}

\begin{proof}
For $s\le y\le Y$ let $\mathbf u(y):=(u_1(y),\dots,u_{2s}(y))$. These are the
cells $y+\lfloor c/2\rfloor\le n$ at the times $2n-2+\lceil c/2\rceil-y$ with
$\lceil c/2\rceil\le s\le y$, so they are defined, and none of them is $\Fq$ when
$y>s$. By \eqref{eq:slices}, $u_c(y)$ with $1\le c\le 2s$ is computed from
components of index less than $c$ of $\mathbf u(y)$, components of index at
most $c$ of $\mathbf u(y+1)$, and the input cells $u_{-1}(y)$ and $u_0(y)$.
Hence $\mathbf u(y)=\Gamma(\mathbf u(y+1),u_{-1}(y),u_0(y))$ for $s+1\le y\le
Y-1$, with $\Gamma$ depending on $\delta$ only. Unfolding \eqref{eq:slices} for
even $c$, $u_{2j}(y)=\delta(u_{2j-2}(y),u_{2j-1}(y+1),u_{2j}(y+1))$, shows that
$u_{2s}(y)=\Gamma'(\mathbf u(y+1),u_0(y))$ for $s\le y\le Y-1$. By hypothesis,
$u_0$ is $\ell$-periodic on $[s,Y]$ and $u_{-1}$ on $[s+1,Y]$.

Suppose that $Y-s>(k-1)^{2s}\ell$. The pairs $(\mathbf u(y),y\bmod\ell)$ with
$s+1\le y\le Y$ take at most $(k-1)^{2s}\ell$ values, so there are $s+1\le
y_1<y_2\le Y$ with $\mathbf u(y_1)=\mathbf u(y_2)$ and $\ell\mid h:=y_2-y_1$.
By downward induction, $\mathbf u(y)=\mathbf u(y+h)$ for $s+1\le y\le y_1$: if
this holds for $y+1\le y_1$, then
\[
\mathbf u(y)=\Gamma(\mathbf u(y+1),u_{-1}(y),u_0(y))
=\Gamma(\mathbf u(y+1+h),u_{-1}(y+h),u_0(y+h))=\mathbf u(y+h),
\]
since $y+h<y_2\le Y$. As $s+h\le Y-1$ and $u_0(s+h)=u_0(s)$,
\[
u_{2s}(s+h)=\Gamma'(\mathbf u(s+1+h),u_0(s+h))=\Gamma'(\mathbf
u(s+1),u_0(s))=u_{2s}(s)=\Fq,
\]
although $\lceil 2s/2\rceil=s<s+h$, so that $u_{2s}(s+h)$, the cell $2s+h$ at
time $2n-2-h$, is not $\Fq$. This is a contradiction.
\end{proof}

In words: if the two input anti-diagonals are $\ell$-periodic over a stretch
of more than $(k-1)^{2s}\ell$ cells that starts at the cell $s$, then the band
of the first $2s$ chains repeats a state, and the cell $2s+h$ fires $h$ steps
too early (\cref{fig:leftproof}). The bound is exponential in $s$, but for $s=1$ and $s=2$ it is small
enough to be useful in computations (\cref{prop:germs,cor:germ12}).

\begin{figure}[t]
\centering
%% Figure: the proof of the left-border lemma (Lemma 4.1). Schematic: the lower left corner of the
%% line of length n = 20 with s = 2, l = 2, Y = 16, y_1 = 8, y_2 = 12, h = 4. The cell (t,i), 20 <= t <= 38,
%% is [i-1,i] x [19-t,20-t]; chain c is the anti-diagonal t+i = 2n-2+c = 38+c.
\begin{tikzpicture}[x=2.6mm,y=2.6mm]
  \foreach \t in {20,...,38} {\foreach \i in {1,...,18} {%
    \pgfmathtruncatemacro{\c}{\t+\i-38}%
    \pgfmathtruncatemacro{\cls}{(\t==38) ? 1 : ((\c<=-2) ? 2 : ((\c==-1) ? (((\i>=2)&&(\i<=15)) ? 3+mod(\i,2) : 5)
      : ((\c==0) ? (((\i>=2)&&(\i<=16)) ? 3+mod(\i,2) : 5) : ((\c<=4) ? 6+mod(\c,2) : 8))))}%
    \def\cc{\ifcase\cls\or stF\or rgAgree\or stB\or stA\or rgInput\or rgRefl\or rgRefl2\or gNone\fi}%
    \filldraw[fill=\cc,draw=grid!45,line width=0.15pt] (\i-1,20-\t) rectangle (\i,19-\t);}}
  \draw[ink,line width=1.2pt] (0,0.2) -- (0,-19.2);
  % slices y_1 = 8 and y_2 = 12 (equal), slice s = 2 and the cell u_{2s}(s+h)
  \def\slice#1#2{\draw[#2] (#1-1,1-39+#1+19) -- ++(2,0) -- ++(0,-1) -- ++(1,0) -- ++(0,-1) -- ++(-2,0) -- ++(0,1) -- ++(-1,0) -- cycle;}
  \slice{8}{ink,line width=1.3pt}
  \slice{12}{ink,line width=1.3pt}
  \slice{2}{ink,line width=1.3pt,densely dashed}
  \fill[stF] (7,-14) rectangle (8,-15);
  \fill[pattern=north east lines,pattern color=white] (7,-14) rectangle (8,-15);
  \draw[ink,line width=1.3pt] (7,-14) rectangle (8,-15);
  % labels below the input anti-diagonals and of the chains
  \node[font=\scriptsize,anchor=south] at (15.5,-0.0) {$Y$};
  \node[font=\scriptsize,anchor=east] at (-0.3,-0.5) {$t=n$};
  \node[font=\scriptsize,anchor=east] at (-0.3,-18.5) {$t=2n-2$};
  \foreach \i/\l in {1/1,2/2,4/4,8/8,18/18} \node[font=\scriptsize,anchor=north] at (\i-0.5,-19.3) {$\l$};
  % annotations
  \draw[-{Stealth[length=1.5mm]},ink,line width=0.5pt] (20.5,-1.5) -- (13.7,-4.0);
  \node[font=\scriptsize,anchor=west,align=left] at (20.5,-1.5) {input anti-diagonals $t+i=2n-3,\,2n-2$:\\
    $\ell$-periodic from the cell $s$ to the cell $Y$};
  \draw[-{Stealth[length=1.5mm]},ink,line width=0.5pt] (20.5,-6.6) -- (14.3,-8.9);
  \draw[-{Stealth[length=1.5mm]},ink,line width=0.5pt] (20.5,-6.6) -- (10.3,-12.9);
  \node[font=\scriptsize,anchor=west,align=left] at (20.5,-6.6) {equal slices $\mathbf u(y_1)=\mathbf u(y_2)$, with\\
    $y_1=8$, $y_2=12$ and $\ell\mid h=y_2-y_1$};
  \draw[-{Stealth[length=1.5mm]},ink,line width=0.5pt] (20.5,-12.6) -- (8.3,-14.4);
  \node[font=\scriptsize,anchor=west,align=left] at (20.5,-12.6) {then $u_{2s}(s+h)=u_{2s}(s)=\Fq$:\\
    the cell $2s+h$ fires at time $2n-2-h$};
  \draw[-{Stealth[length=1.5mm]},ink,line width=0.5pt] (20.5,-17.4) -- (4.3,-18.6);
  \node[font=\scriptsize,anchor=west,align=left] at (20.5,-17.4) {slice $s$ (dashed): $u_{2s}(s)$ is\\
    the cell $2s$ of the firing row};
  % legend
  \node[anchor=north west,font=\scriptsize,align=left] at (0,-21.6) {%
    \swatch{stB}\,\swatch{stA}\,periodic part of the input, shaded by $i\bmod\ell$\quad
    \swatch{rgRefl2}\,\swatch{rgRefl}\,chains $1,\dots,2s$\quad \swatch{gNone}\,chains $c>2s$\quad \swatch{stF}\,$\Fq$};
\end{tikzpicture}

\caption{The proof of \cref{lem:left}, drawn for $n=20$, $s=2$, $\ell=2$ and
$Y=16$: the lower left corner of the line of length $n$, at the times
$n\le t\le2n-2$. The two input anti-diagonals $t+i=2n-3$ and $t+i=2n-2$ are
$\ell$-periodic from the cell $s$ to the cell $Y$. Each slice $\mathbf u(y)$,
a group of $2s$ cells on the chains $1,\dots,2s$ (outlined), is computed from
the slice $\mathbf u(y+1)$ and from two input cells, so the slices form a
finite automaton that reads the input anti-diagonals from right to left. If two
slices are equal, $\mathbf u(y_1)=\mathbf u(y_2)$ with $\ell\mid
h=y_2-y_1$, then $\mathbf u(y)=\mathbf u(y+h)$ for $s+1\le y\le y_1$, and the
last cell $u_{2s}(s+h)$ of the slice $s+h$ (hatched) has the state of the last
cell $u_{2s}(s)$ of the slice $s$, which is $\Fq$. It is the cell $2s+h$ at
the time $2n-2-h$, which would fire $h$ steps too early. Hence $Y-s\le
(k-1)^{2s}\ell$.}
\label{fig:leftproof}
\end{figure}

\subsection{No Periodic Wedge at the Left Border}

\begin{theorem}[left-border barrier]\label{thm:left}
Let $\delta$ be a minimal-time solution. There are no $s\ge1$, $\alpha>0$,
$t_0$ and $p,q\ge1$ such that $C_\infty(t+p,i)=C_\infty(t,i)$ and
$C_\infty(t,i+q)=C_\infty(t,i)$ whenever the cells involved lie in the wedge
\[
\Omega:=\{(t,i):\ s\le i\le\alpha t,\ t\ge t_0\}.
\]
\end{theorem}

\begin{proof}
Suppose otherwise and put $\ell:=pq$. If $(t,i)$ and $(t-\ell,i+\ell)$ lie in
$\Omega$, so do the cells $(t-\ell,i+jq)$ for $0\le j\le p$ and $(t-\ell+jp,i)$
for $0\le j\le q$, hence
$C_\infty(t-\ell,i+\ell)=C_\infty(t-\ell,i)=C_\infty(t,i)$. Let
$Y_n:=\lfloor\alpha(2n-3)/(1+\alpha)\rfloor$. For $s\le y\le Y_n$ the cells
$(2n-2-y,y)$ and $(2n-3-y,y)$ satisfy $y\le\alpha(2n-3-y)$, so they lie in
$\Omega$ as soon as $2n-3-Y_n\ge t_0$, that is, for all large $n$. A step of
$\ell$ positions along an anti-diagonal is the step from $(t,i)$ to
$(t-\ell,i+\ell)$, so the words of \cref{lem:left} are $\ell$-periodic on
$[s,Y_n]$ and $[s,Y_n-1]$. As $Y_n\ge2\alpha n/(1+\alpha)-4$, for large
$n$ the number $Y:=\min(Y_n,n-s)$ satisfies $Y-s>(k-1)^{2s}\ell$, which
contradicts \cref{lem:left} (\cref{fig:wedge}(a)).
\end{proof}

\begin{figure}[t]
\centering
%% Figure: the left-border barrier (Theorem 4.2) and eventually regular half-lines (Corollary 4.3).
%% Schematic, in continuous coordinates: position i to the right, time t downwards.
\begin{tikzpicture}[x=1.6mm,y=1.6mm]
  % (a) the wedge and the input anti-diagonals of three lengths (s = 2, alpha = 0.4, t_0 = 8)
  \begin{scope}
    \fill[gNone] (0,0) -- (40,-40) -- (0,-40) -- cycle;
    \fill[stB] (2,-8) -- (3.2,-8) -- (16,-40) -- (2,-40) -- cycle;
    \draw[stC,line width=0.6pt] (2,-40) -- (2,-8) -- (3.2,-8) -- (16,-40);
    \foreach \n in {12,16,20} {
      \draw[grid,line width=0.5pt] (\n-1,-\n+1) -- (0,-2*\n+2);
      \pgfmathsetmacro{\Y}{0.4*(2*\n-3)/1.4}
      \draw[rgCone,line width=2.2pt] (2,-2*\n+4) -- (\Y,-2*\n+2+\Y);}
    \draw[stF,line width=0.9pt] (5,-25) -- (5,-22) -- (8,-22);
    \fill[ink] (5,-25) circle (0.45) (8,-22) circle (0.45);
    \draw[ink,line width=1.2pt] (0,0) -- (0,-40);
    \draw[ink,line width=0.7pt] (0,0) -- (40,-40);
    \node[font=\small] at (5.6,-35) {$\Omega$};
    \node[font=\scriptsize,anchor=north] at (2,-40.2) {$s$};
    \node[font=\scriptsize,anchor=north] at (16,-40.2) {$i=\alpha t$};
    \node[font=\scriptsize,anchor=east] at (-0.3,-8) {$t_0$};
    \node[font=\scriptsize,rotate=-45,anchor=south] at (30,-30) {front};
    \node[font=\scriptsize,anchor=west] at (8.4,-21.4) {$(t-\ell,i+\ell)$};
    \node[font=\scriptsize,anchor=north west] at (5.2,-25.1) {$(t,i)$};
    \node[font=\scriptsize,anchor=west,align=left] at (12.4,-11.6) {$t+i=2n-2$};
    \draw[-{Stealth[length=1.4mm]},ink,line width=0.45pt] (12.6,-12.3) -- (6.8,-15.0);
    \node[font=\small,anchor=south west] at (-2,0.8) {(a) the proof of \cref{thm:left}};
  \end{scope}
  % (b) an eventually regular half-line: blocks separated by signals of speeds 0, 1/4 and 1/2
  \begin{scope}[shift={(48,0)}]
    \fill[gNone] (0,0) -- (40,-40) -- (0,-40) -- cycle;
    \fill[rgInput] (0,-8) rectangle (3,-40);
    \fill[stB] (3,-8) -- (11,-40) -- (3,-40) -- cycle;
    \fill[rgRefl2] (3,-8) -- (4,-8) -- (20,-40) -- (11,-40) -- cycle;
    \fill[rgAgree] (4,-8) -- (8,-8) -- (40,-40) -- (20,-40) -- cycle;
    \draw[ink,line width=0.6pt] (3,-8) -- (3,-40) (3,-8) -- (11,-40) (4,-8) -- (20,-40);
    \draw[ink,densely dashed,line width=0.9pt] (4,-20) -- (8,-40) (4,-20) -- (4,-40);
    \draw[grid,line width=0.4pt] (0,-8) -- (8,-8);
    \draw[ink,line width=1.2pt] (0,0) -- (0,-40);
    \draw[ink,line width=0.7pt] (0,0) -- (40,-40);
    \node[font=\scriptsize,anchor=north] at (1.5,-40.2) {$w_0$};
    \node[font=\scriptsize,anchor=north] at (7.0,-40.2) {$z_1^{a_1+\nu D_1}$};
    \node[font=\scriptsize,anchor=north] at (15.5,-40.2) {$z_2^{a_2+\nu D_2}$};
    \node[font=\scriptsize,anchor=north] at (30.0,-40.2) {$z_3^{a_3+\nu D_3}$};
    \node[font=\scriptsize,anchor=east] at (-0.3,-8) {$t_1$};
    \node[font=\scriptsize,anchor=west] at (10.5,-2.5) {transient part};
    \draw[-{Stealth[length=1.4mm]},ink,line width=0.45pt] (10.6,-3.0) -- (3.0,-5.4);
    \node[font=\small,anchor=south west] at (-2,0.8) {(b) an eventually regular half-line};
  \end{scope}
  % legend
  \node[anchor=north west,font=\scriptsize,align=left] at (0,-44.6) {%
    \swatch{stB}\,in (a): the wedge $\Omega$; in (b): the first block that grows, with a wedge (dashed)\\[0.3ex]
    \swatch{rgCone}\,in (a): the stretch from the cell $s$ to the cell $Y_n$ of the input anti-diagonals
    of the lengths $n=12,16,20$};
\end{tikzpicture}

\caption{(a)~The proof of \cref{thm:left}. If the half-line were periodic in
time and in space on the wedge $\Omega=\{(t,i):s\le i\le\alpha t,\ t\ge t_0\}$,
then the input anti-diagonals $t+i=2n-3$ and $t+i=2n-2$ of the line of length
$n$ would be $\ell$-periodic, with $\ell=pq$, from the cell $s$ to the cell
$Y_n\approx2\alpha n/(1+\alpha)$ (thick), a stretch whose length grows
linearly with $n$; for large $n$, \cref{lem:left} gives a contradiction. The
path from $(t,i)$ to $(t-\ell,i+\ell)$ shows why: $\ell$ steps along an
anti-diagonal are $q$ time periods $p$ followed by $p$ space periods $q$
inside $\Omega$. (b)~An eventually regular half-line (\cref{cor:regular}):
from the time $t_1$ on, it consists of finitely many periodic blocks
separated by signals of rational speeds, here $0$, $1/4$ and $1/2$. The first
block that grows contains a wedge (dashed) on which the half-line is periodic
in time and space, which (a) excludes.}
\label{fig:wedge}
\end{figure}

The half-line of a rule is \emph{eventually regular} if there are $t_1$,
$P\ge1$ and, for each residue $r\in\{0,\dots,P-1\}$, an integer $m$, words
$w_0,\dots,w_m$, nonempty words $z_1,\dots,z_m$ over $W$ and integers
$a_j,D_j\ge0$ (all depending on $r$) such that for all $\nu\ge0$ the
configuration $C_\infty(t,1)\cdots C_\infty(t,t+1)$ at the time
$t=t_1+r+\nu P$ is $w_0z_1^{a_1+\nu D_1}w_1\cdots z_m^{a_m+\nu D_m}w_m$. In
words: from some time on, the half-line consists of finitely many periodic
regions separated by signals of rational speeds; on signals in one-dimensional
cellular automata see~\cite{MazoyerTerrier99}. In simulations, the half-line of
the five-state rule $\delta_{12}$ of \cref{sec:five} looks like this
(\cref{fig:halflines}): the cells between roughly the cell $t/3$ and the front
are in state $\Gq$, and the cells between the left border and roughly the cell
$t/3$ alternate between $\Aq$ and $\Lq$, both in time and in space.

\begin{corollary}\label{cor:regular}
The half-line of a minimal-time solution is not eventually regular.
\end{corollary}

\begin{proof}
Suppose that it is. For each residue $r$ some $D_j$ is positive, because the
configuration grows by $P$ cells per period; let $j(r)$ be the least such
index. The blocks and the junctions to the left of $z_{j(r)}$ have constant
lengths, so the block $z_{j(r)}^{a_{j(r)}+\nu D_{j(r)}}$ starts at a fixed cell
$s_r$, and its length is at least $\beta t$ for some $\beta>0$ and all large
$t$ of the residue class $r$. Inside this block, the state of a cell depends
only on its distance to $s_r$ modulo $|z_{j(r)}|$, and it is the same at the
times $t$ and $t+P$. Let $s:=\max_r s_r$, $q:=\prod_r|z_{j(r)}|$ and
$\alpha:=\beta/2$. For $t_0$ large, every cell of the wedge $\{(t,i):s\le
i\le\alpha t,\ t\ge t_0\}$ lies in the block of index $j(r)$ of its residue
class, so $C_\infty$ is periodic with period $P$ in time and $q$ in space on
this wedge (\cref{fig:wedge}(b)). This contradicts \cref{thm:left}.
\end{proof}
